\chapter{Triple Integrals}

How much does a scoop of ice cream weigh if it is denser at the bottom where it
has started to melt? How much air is inside a dome-shaped tent? Questions about
solids need one more dimension of adding up. That is what triple integrals do.

\section{From Double to Triple Integrals}\index{triple integral}
A double integral adds up a function over a flat region. A
\index{triple integral}\newterm{triple integral} adds up a function over a
three-dimensional solid $E$:
$$\iiint_E f(x, y, z)\,dV$$
If $f = 1$, the triple integral gives the volume of $E$. If $f$ is a density,
it gives the mass of $E$.

\section{Iterated Triple Integrals}\index{iterated integral!triple}
We evaluate a triple integral one variable at a time, from the inside out, just
as we did for double integrals. For a rectangular box
$0 \leq x \leq 1$, $0 \leq y \leq 2$, $0 \leq z \leq 3$:
$$\int_0^1 \int_0^2 \int_0^3 xyz\,dz\,dy\,dx
= \int_0^1 x\,dx \cdot \int_0^2 y\,dy \cdot \int_0^3 z\,dz
= \frac{1}{2} \cdot 2 \cdot \frac{9}{2} = \frac{9}{2}$$
Because the limits are constants, the integral splits into a product of three
one-variable integrals.

When the limits depend on other variables, the order matters. Consider the
solid $E$ in the first octant bounded by the plane $x + y + z = 1$. Its volume is
$$\int_0^1 \int_0^{1-x} \int_0^{1-x-y} dz\,dy\,dx$$
The innermost integral gives $1 - x - y$. The middle integral gives
$\frac{(1-x)^2}{2}$. The outer integral gives
$\int_0^1 \frac{(1-x)^2}{2}\,dx = \frac{1}{6}$. The volume is $\frac{1}{6}$.

\section{Cylindrical Coordinates}\index{cylindrical coordinates!triple integral}
Many solids are round around an axis. For those, use cylindrical coordinates,
where $x = r\cos\theta$, $y = r\sin\theta$, and $z$ stays the same. The volume
element becomes $dV = r\,dz\,dr\,d\theta$. The extra factor of $r$ comes from
the area of a thin ring, $2\pi r\,dr$.

For a cylinder of radius 2 and height 3:
$$\int_0^{2\pi} \int_0^2 \int_0^3 r\,dz\,dr\,d\theta
= 2\pi \cdot \int_0^2 3r\,dr = 2\pi \cdot 6 = 12\pi$$
This matches the formula $\pi r^2 h = \pi \cdot 4 \cdot 3 = 12\pi$. Think of a can of soup: its volume is the area of the round top times the height. The triple integral gives the same answer.

\textbf{Example}: Find the volume of an upside-down cone whose tip is at the
origin and whose top is the disk of radius 1 at height 1.

\textbf{Solution}: The cone's side is $z = r$. For each point $(r, \theta)$ in
the unit disk, $z$ runs from the side up to the top, $r \leq z \leq 1$:
$$\int_0^{2\pi}\int_0^1\int_r^1 r\,dz\,dr\,d\theta = 2\pi\int_0^1 r(1 - r)\,dr
= 2\pi\left(\frac{1}{2} - \frac{1}{3}\right) = \frac{\pi}{3}$$
This matches $\frac{1}{3}\pi r^2 h$ with $r = h = 1$.

\section{Spherical Coordinates}\index{spherical coordinates!triple integral}
For solids that are round in every direction, use spherical coordinates. Let
$\rho$ be the distance from the origin, $\varphi$ the angle down from the positive
$z$-axis, and $\theta$ the angle around the $z$-axis as before:
$$x = \rho\sin\varphi\cos\theta, \quad y = \rho\sin\varphi\sin\theta, \quad
z = \rho\cos\varphi$$
The volume element is $dV = \rho^2\sin\varphi\,d\rho\,d\varphi\,d\theta$.

For the unit ball:
$$\int_0^{2\pi} \int_0^{\pi} \int_0^1 \rho^2\sin\varphi\,d\rho\,d\varphi\,d\theta
= 2\pi \cdot 2 \cdot \frac{1}{3} = \frac{4\pi}{3}$$
This is the familiar volume formula for a sphere, $\frac{4}{3}\pi r^3$, with $r = 1$.

\textbf{Example}: A ball of radius 1 is denser near its center, with density
$\delta = 2 - \rho$ (in g/cm$^3$, with $\rho$ in cm). Find its mass.

\textbf{Solution}:
$$m = \int_0^{2\pi}\int_0^{\pi}\int_0^1 (2 - \rho)\rho^2\sin\varphi\,d\rho\,d\varphi\,d\theta
= 2\pi \cdot 2 \cdot \left(\frac{2}{3} - \frac{1}{4}\right) = \frac{5\pi}{3} \approx 5.24 \text{ g}$$
A ball with uniform density 1 would have mass $\frac{4\pi}{3} \approx 4.19$ g, so
the extra density near the center adds about a gram.

\section{Choosing a Coordinate System}
Use rectangular coordinates when the limits are flat planes. Use cylindrical
coordinates when the solid is round around one axis, such as a cylinder, a
cone, or a paraboloid $z = x^2 + y^2$. Use spherical coordinates when the solid
is round around a point, such as a ball or a sphere.

\begin{Exercise}[label=triple1, title=Evaluate a box integral]
Evaluate $\displaystyle \int_0^1 \int_0^1 \int_0^1 (x + y + z)\,dz\,dy\,dx$.
\end{Exercise}
\begin{Answer}[ref=triple1]
Each of the three terms contributes $\frac{1}{2}$, so the total is $\frac{3}{2}$.
\end{Answer}

\begin{Exercise}[label=triple2, title=Volume by cylindrical coordinates]
Use cylindrical coordinates to find the volume of the solid inside the cylinder
$r = 1$, between $z = 0$ and $z = 4$.
\end{Exercise}
\begin{Answer}[ref=triple2]
$\int_0^{2\pi}\int_0^1\int_0^4 r\,dz\,dr\,d\theta = 2\pi \cdot \frac{1}{2} \cdot 4 = 4\pi$.
\end{Answer}

\begin{Exercise}[label=triple3, title=Volume of a larger ball]
Use spherical coordinates to find the volume of a ball of radius 2.
\end{Exercise}
\begin{Answer}[ref=triple3]
$\int_0^{2\pi}\int_0^{\pi}\int_0^2 \rho^2\sin\varphi\,d\rho\,d\varphi\,d\theta
= 2\pi \cdot 2 \cdot \frac{8}{3} = \frac{32\pi}{3}$.
\end{Answer}

\begin{Exercise}[label=triple4, title=Order of integration]
Write the volume of the solid in the first octant under $x + y + z = 1$ as an
iterated integral in the order $dx\,dy\,dz$.
\end{Exercise}
\begin{Answer}[ref=triple4]
$\displaystyle \int_0^1 \int_0^{1-z} \int_0^{1-y-z} dx\,dy\,dz = \frac{1}{6}$.
\end{Answer}
